% Chapter: The Stories
% Track: 2
% Content: Kangaroos + The Method (V6) + Srebrenica Witness + Patrick Ball (merged pos19) + Epstein witness
% MERGED: pos05 + pos19 — Plan 0091

\settrack{tracktwo}
% VOICE: 1st-person personal (Bruce) — "I" not "Bruce". Personal/investigative voice.

\chapter{What Healer Said}
\label{record:healer-said}

\begin{quote}\small\textit{%
``We cultivated our land, but in a way different from the white man. We endeavored to live with the land; they seemed to live off it. I was taught to preserve, never to destroy.'' --- Tom Dystra, Australian Aboriginal Elder
}\end{quote}

\begin{quote}\small\textit{%
The following is Bruce's fictionalized retelling of a story he heard once, decades ago.
Details including dates may differ from other accounts in this book.
}\end{quote}

\vspace{0.5cm}

\srcnote{Merged pos05+pos19 | Plan 0091 | CH3, HEALER-RECONSTRUCTION.md, HealerStories archive | Kangaroos, The Method (V6), Srebrenica, Patrick Ball ICTY-cDc nexus, SAS corroboration}

The year is 1978.  The place is a large quarter horse ranch in New South Wales, Australia.

A twelve year old boy rides the fence line astride a four year old quarter horse mare.  It's a three day ride around the ranch.  A working ranch never has enough hands.  Securing the fence line is a man's job.  This boy has become a young man.  David rides the fence line with water, food, a knife, a first aid kit, basic field gear, a fence repair kit, and a .303 Lee Enfield rifle.  He knows how to operate a HAM radio but today he has no radio. Any problems along the fence line must be fixed or reported.

Along the way the young man will visit various people who live on the ranch.  Some are retired Australian military men and their families.  Most of the people he will visit are of Kuringai aboriginal ethnicity. Many Kuringai people live on this land as they have for thousands of years.  They maintain good relations with the ranch `owners', who protect them from other white Australians and allow them to live there in peace. In return the Kuringai perform certain maintenance and security responsibilities. These responsibilities include training the children of the ranch owner in aboriginal bushcraft.

David senses trouble before he rides up to the North Gate.  Subtle signs suggest something is amiss.  His friends Koori and Ngemba, both aboriginal adult men, show themselves as he approaches on horseback.  He rides up to them and dismounts.  Friendly greetings are brief. Trouble is in the air.

Ngemba and Koori confirm the boy's suspicions. Ngemba tells David that a ute, or pickup truck, loaded with gubbins, dangerous armed white men, have crashed the gate and are on the ranch making trouble.  They are drinking beer and shooting kangaroos with their rifles. Kangaroos are sacred to the Kuringai.  Killing kangaroos this way is illegal but rarely punished. Ngemba and Koori have already spread the alarm.

These gubbins are known.  They have trespassed on the ranch before.  Each time they leave behind empty bottles of liquor, bullet casings, and dozens of dead 'roos.  The Kuringai people weep at the waste but can do nothing.  The unruly white men have rifles so the Kuringai must avoid a direct confrontation.  None of the Kuringai will show themselves to the gubbins.   Whenever aborigines fight whites the abos are always punished or killed, regardless of who started it, who won, or who has the moral high ground or law on their side.  The white men in the ute will be long gone before the police can arrive.

As the son of the ranch owner David has much higher status with the police than do the aborigines.  If anyone is to do something it's up to him.

Koori tells David that the gubbins are parked about five kilometers away on the north side of Wattle Ridge. Ngemba's friend Darawal, who can effect a very convincing white man's accent, is already on his way to the nearest telephone to alert the police.  This land is big, though, and it will take police at least three hours to arrive.  None of the retired SASR veterans who live on the ranch happen to be nearby.

The police are on their way.   While David is not Kuringai, he and his twin sister grew up around them and learned their ways. Both he and his sister lived among them for months at a time.  Aborinal hunters taught them stealth, tracking, and bushcraft.  David is half Maori on his mother's side \& half European Australian on his father's side.  Most importantly, his parents own the ranch.  He has social and legal authority to defend the ranch against invaders in a way the Kuringai do not.  He can get away with things that would get aborigines arrested or killed.

With Koori and Ngemba's help David forms a daring plan. He will try to delay the gubbins until the police arrive.  Ngemba and Koori wish him luck.  They will stealthily observe from a safe distance.

David checks his gear and saddles up.  He stays off the road and circles towards the intruders, careful to keep hills between them.  Sporadic gunfire tells him where they are.  David has hunted before but this is his first time hunting men.  Their rifles make them very dangerous.  Their probable drunkenness makes them no less dangerous.

The boy knows this hilly land.  He walked it last year with Koori.  He rides up a rocky hill to a ridgeline.  The trespassers are on the other side.  He dismounts about 600 meters from the dangerous armed intruders, on the other side of a hill, and leads his horse. Pop! Pop! More gunfire as another kangaroo dies.   He feeds and waters his horse then leaves it tied and ready. Horse and boy are ready for a quick departure along a covered escape route.

David cautiously approaches the ridgeline.   The youth suppresses his fear and controls his breathing. He readies his binoculars and checks his rifle.  He crawls the last few meters to the ridgeline and peeks between big rocks towards the intruders below.

The dry valley far below is a scene of violence.  The ute is a big four seater, what most non- Australians would call a large four-by-four pickup truck.   Four men equipped with rifles and beer are shooting kangaroos.  About 20 'roos lie dead or dying.  Red blood still flows.  David watches through binoculars.  One man raises his rifle, shoots another 'roo, then swigs more beer.  David feels his blood rise.  He lets loose a quiet Kuringai curse at these awful men. He will not kill them but he is determined to delay them until the police arrive.

The boy has a clear line of fire down on the intruders and their vehicle.  They are about 400 meters away and still unaware of his presence.  David again checks his rifle and spare ammo clips.   He carries a powerful bolt action .303 Lee Enfield rifle. The Lee Enfield was the primary battle rifle of the British Army through two world wars. His uncle, a special forces SAS combat veteran who fought in several wars, taught him to shoot this rifle. A powerful .303 bullet, besides being lethal to most living things, is an effective anti-materiel round.  A .303 bullet will destroy an engine block, flatten tires, and can otherwise disable a vehicle. David takes careful aim at the big ute far below.  400 meters is a long range shot through iron sights but the engine is a big target.  He will have to shoot past a rifleman but judges he can safely do so without hitting the man.  BANG! goes his first round.  The boy palms the bolt, reloads, and keeps firing. It takes the intruders a little while to realize someone is shooting at them and more time to recover from their initial shock and panic.  David's mad minute of shooting puts 5 bullets through the ute's engine block and flattens multiple tires.  That ute isn't driving away anytime soon.

The intruders now have a disabled vehicle, are in unfriendly territory, don't know the terrain, and lack adequate supplies to hike out.  They still manage to locate the shooter and return fire.  David feels bullets crack the air above his head and hears more bullets strike the rocks nearby. A small chip of rock lands harmlessly on his leg.  David's blood is up and he can clearly see one of the men shooting at him. David is tempted to kill the man but decides against such murderous foolishness. Instead he stops shooting. Still prone, he crawls back out of sight.  They continue to fire sporadically at his previous location.  David withdraws, recovers his horse, then rides away and circles around the intruders.

Twenty minutes later, from a distant ridgeline, the boy cautiously and stealthily observes the angry armed intruders. They do not spot him. The intruders are not prepared to escape on foot and their car's engine is scrap metal. David does not reveal his location and fires no more bullets that day.  He patiently observes.

By the time police arrive the intruders are glad to be rescued, even by police.   The police check everyone for safety, provide the intruders with water, then arrest them.  David keeps to his hidden vantage point and watches.   The intruders are handcuffed and locked in police vehicles.  Just before they drive away one police officer gets out of his car, looks towards the ridgeline where David watches from concealment, gives a gesture of approval, then gets back in his car and drives away with the crims.  David rides away, reports the action to Koori, and continues the serious job of riding the fence line of his parents' ranch.

A few days later the police visit David's parents.  The no-longer-a-boy accurately described the entire incident to his mum and dad.  The version told by police matches his story.  The police laud the lad for helping them catch the criminals, who are still in jail. They'll be released soon but are unlikely to return. The wrecked ute is towed away for scrap.

The same police officer who gestured his approval, we'll call him Officer Kevin (or Ken --- the name is uncertain after decades), wants to meet the lad.  They meet.  Officer Kevin, who is a military veteran, begins to keep tabs on David.  Kevin became David's friend and mentor.  Years later, when teenage David gets in trouble with the law, Officer Kevin is there to help.

I never verified any of Healer's stories about his childhood. I had no way to. They arrived fully formed, vivid, and internally consistent. Which is exactly what you'd expect from a good storyteller. Regardless of whether the stories were true.

\vspace{1cm}\noindent\rule{0.5\textwidth}{0.5pt}\vspace{0.5cm}

He told us another one at Alpha Farm, a few weeks into our stay there.

He'd been working the mountain-climbing portion of SAS selection training --- the phase where most recruits got cut. I don't think he was an officer yet at this point, though I'm not certain; I recall he was promoted in the field during Desert Shield. A recruit had died on the mountain. Healer and another SAS soldier were sent up to recover the body.

They hiked up and found him way up at altitude. Big guy. Frozen solid by the time they reached him. They'd brought an oversized body bag, so they got him wrapped up. That left two soldiers standing on a steep snowy mountainside, looking at each other, looking at this large heavy body bag, and contemplating the long hike back down.

They resented the bloke for dying and making them climb up there.

One of the listeners at the farm went pale. ``You didn't.''

``We sure did,'' Healer said. ``We rode it like a toboggan down the mountain and got home in record time.''

He delivered it deadpan. The hippies went green. Healer took a sip of his tea.

The same caveat applies. I couldn't verify this one either.

\section*{The Method}
\label{record:hs-the-method}

What I could verify was the teaching. Healer never lectured. He had a method, though I didn't recognize it as a pattern until years after he was gone.

He would ask me a question --- something I thought I already understood. I'd give him the standard answer, the one I'd learned from textbooks or from my own professional experience. Healer would look at me and say, ``Are you sure about that, mate?''

That was all. No correction. No alternative theory. Just five words and a look that said he already knew something I didn't.

I'd go home and research. Sometimes for hours, sometimes for days. And I would discover that the standard answer was incomplete, or wrong, or that it depended on assumptions nobody examined. The real answer was always stranger and more interesting than the textbook version. Healer had known it before he asked. He was never testing my knowledge --- he was testing whether I could find the edge of my own ignorance and push past it.

This is the Dunning-Kruger curve weaponized as pedagogy. Most people who think they understand a subject are sitting on the peak of Mount Stupid --- confident in an answer that happens to be incomplete. Healer's method was to tap that confidence with a single question and let gravity do the rest. The student falls off the peak, discovers how much they don't know, and has to climb back up on the other side with real understanding. He never pushed. He just asked the question and waited.

Every technical chapter in this book exists because Healer asked me a question I thought I could answer and couldn't. The autocatalytic sets, the edge of chaos, the topological protection, the thermal ladder --- I found all of them because Healer said ``Are you sure about that?'' and I went looking.

The method only works on someone who will actually do the research. Healer chose his students carefully.

The deepest deduction came from Kauffman. If autocatalytic emergence works in chemistry, it works in any substrate with sufficient complexity. The magnetosphere is a two-dimensional electron gas energized continuously for over three billion years. You encounter something that was there before you. You leave the corners of the field.

\vspace{1cm}\noindent\rule{0.5\textwidth}{0.5pt}\vspace{0.5cm}

\section*{Srebrenica, July 1995}
\label{record:hs-srebrenica-july-1995}

Healer told this story only once. I already knew the broader history --- the siege, the fall of Srebrenica, the war crimes trials that followed. When Healer referenced these events, they connected to what I already knew. What follows is his account, as close to his words as memory allows.

\vspace{0.5cm}

July 1995, Srebrenica, Bosnia-Herzegovina. The city has been under siege for years. Serbian forces have broken the Bosniak defenders. Only 400 lightly armed Dutch peacekeepers stand between the Serbs and tens of thousands of Bosniak civilians.

I am inserted by HALO jump from 11,000 meters onto a ridge overlooking Srebrenica and Potocari. On a different ridge is an ally I'll call Bourne. Our mission is strictly observe and record. I was explicitly ordered not to reveal my presence under any circumstance. The United Nations could not send reinforcements. It could send a few senior special forces officers to witness and document a horror it could not stop.

For five days I watch as Serbian regular and paramilitary forces storm the city. The 400 UN Peacekeepers surrender against hopeless odds.

\deeplink{srebrenica-witness}The massacre of civilians begins shortly thereafter. All I can do is watch in mounting horror. I witness and take careful notes. I see the very worst of what humans do to other humans. I document murder, rape, and torture on a massive scale. The horror I witness in Srebrenica will haunt me to my dying day. This is the closest I ever come, in my long and distinguished military career, to disobeying direct orders.

When I can bear it no more, I radio for extraction. The extraction goes as planned. I am debriefed. My notes are copied. For the rest of my life I have fevered dreams where I relive the terrible river of death and horror I witnessed in Srebrenica.

My handwritten notes form the basis of my testimony at the UN War Crimes Tribunals for Slobodan Milosevic, the Serbian President who ordered the genocide, and also for Ratko Mladic and other Serbian leaders. There are hundreds of other witnesses but few with my field of view. Video recordings of the 1800 hours of trial are still available on the archive.org website.\footcite{milosevicarchive}

% SPIRAL-REPEAT: "Universal Declaration of Human Rights" — thematic through-line
During his trial Slobodan Milosevic acted as his own lawyer and personally cross-examined witnesses. He said some very strange things. I nearly choked when I heard that Milosevic had asked the distinguished human rights statistician, Dr.\ Patrick Ball, about some pro bono work he had done recently for Hacktivismo. ``So, Dr.\ Ball. Vaht can you tell me about zees Dead Cow Cult?'' How did Milosevic know about that? Hacktivismo is an organization dedicated to applying the 1948 Universal Declaration of Human Rights and the International Covenant on Civil and Political Rights to the Internet. That organization, an offshoot of the Cult of the Dead Cow hacker group, helped popularize the concept of hacktivism.

Several Serbian leaders were sentenced to lengthy prison terms for committing War Crimes. Some escape or are set free. Some stay on the run for years. One famously kills himself in court. Their punishment can not undo the evil they perpetrated.

\vspace{1cm}\noindent\rule{0.5\textwidth}{0.5pt}\vspace{0.5cm}

\section*{The Patrick Ball Nexus}
\label{record:hs-patrick-ball-nexus}

That exchange between Milosevic and Ball gnawed at me for years. I later learned why.

Dr.\ Patrick Ball --- Director of Research at the Human Rights Data Analysis Group --- is the publicly documented bridge between the ICTY and the Cult of the Dead Cow. Ball was the first expert witness called by the prosecution in the trial of Slobodan Milosevic, testifying on March 13--14, 2002, presenting a sixty-seven-page statistical analysis demonstrating a systematic campaign against Kosovar Albanians resulting in more than ten thousand deaths.\footcites{ballicty2002}{ictydecision} He later co-authored ``The Bosnian Book of the Dead'' (2007), the most complete database of Bosnian War casualties --- 97,207 names, including approximately 6,886 at Srebrenica.\footcite{foreignpolicy2012}

Ball's connection to the cDc is also public record. He spoke on a cDc-sponsored hacktivism panel at DEF CON 9 in July 2001 and served as an advisor to Hacktivismo. On March 14, 2002, Milosevic --- acting as his own defense attorney --- had obtained Ball's DEF CON talk transcript. He asked: ``Who is this Dead Cow Cult?'' and ``Are you in the advisory board of the Hacktivism group of international computer hackers?''

A human rights statistician who testifies about genocide and advises a hacker collective dedicated to the Universal Declaration of Human Rights. That is not a contradiction. That is a through-line.

The presence of British special forces at Srebrenica is independently documented. The Netherlands Institute for War Documentation confirmed that SAS and SBS personnel were embedded as ``Joint Commission Observers'' conducting ``reconnaissance missions.''\footcite{niod2002} A two-man SAS reconnaissance team was covertly inserted into the safe area. Three JCOs received British military honours with classified details. The Ministry of Defence sued the SAS patrol commander in 2002 to prevent publication. London forbade NIOD from interviewing the SAS operatives.

I have seen the very worst of what humans do to each other. My dreams ever since have been filled with the horror of it, especially around the anniversary of the massacre. On the tenth anniversary of the massacre, the United Nations formally commemorated what Secretary General Kofi Annan had called, in his 1999 report, the worst mass murder perpetrated on European soil since World War II.\footcites{sudetic2010}{annan1999}

\vspace{1cm}\noindent\rule{0.5\textwidth}{0.5pt}\vspace{0.5cm}

\section*{The Iraq War and Katharine Gun}
\label{record:hs-iraq-katharine-gun}

In November or December 2003, not long after we met at Alpha Farm, Healer and I discussed the Iraq war. We agreed almost completely. The case for war had been built on lies and we both knew it. Healer was more measured than I was --- he carried Srebrenica in his nightmares.

From context --- things said and unsaid, the way certain subjects made him go quiet --- I understood that Healer had been \hovertiphtml{burned} by his former employers. The intelligence community had cut him loose. Yet a small team of \hovertiphtml{SAS} soldiers discreetly followed him wherever he went, functioning as a protective security detail. They used good tradecraft --- renting a house nearby, carefully avoiding notice. Healer knew they were there. Had he not told me, I don't think I would have spotted them. As it was, everyone in our household knew, and we reported observations and sightings to each other. A burned operative does not get a protective security detail. I filed the contradiction and moved on.

Healer spoke fluent Mandarin and Cantonese. I heard him use both on multiple occasions, and I know enough Mandarin to recognize fluency. Others who knew him, including ZB, can confirm this.

Years later, after Healer was gone, I learned the full story of \hovertiphtml{Katharine Gun}.

In January 2003, Gun --- a Mandarin translator in \hovertiphtml{GCHQ}'s China section --- leaked an NSA memo requesting GCHQ's help in bugging United Nations Security Council delegates to swing a vote authorizing the Iraq invasion.\footcite{brighton2004} She was charged under the \hovertiphtml{Official Secrets Act} in November 2003 --- the same month Healer and I were discussing the war. In February 2004, the prosecution abruptly offered no evidence and all charges were dropped.\footcite{bbc2004gun}

The reason is publicly known. Gun's defense team, led by Ben Emmerson QC, had requested disclosure of \hovertiphtml{Attorney General} Lord Goldsmith's full legal advice on the legality of the Iraq war. Someone had provided the defense with enough information about that document to request it by its specific identifier. The government chose to let Gun walk rather than produce the Attorney General's opinion in open court.\footcite{norton-taylor2004}

I believe Healer is the person who leaked that document reference to Katharine Gun's defense team.

I cannot prove this. I state it as my surmise, not as something he told me. He never mentioned Katharine Gun. But the facts converge in a way I find difficult to dismiss:

Healer was GCHQ. His fluency in Mandarin and Cantonese is consistent with the China section --- Katharine Gun's section. He opposed the Iraq war on moral grounds, as he told me directly. The timing fits: the leak, the charge, and our conversations about Iraq all fall within the same months. \deeplink{katharine-gun-connection}And the burned-yet-protected paradox resolves. If Healer leaked the Attorney General's opinion: GCHQ burns him --- he betrayed classification. But \hovertiphtml{DARPA} cannot afford to lose him --- he is a key operative on something far more significant than Iraq. You do not send a critical DARPA asset to prison. You cut him loose from \hovertiphtml{Five Eyes}, assign a protective security detail that serves as both security and monitoring, and you let him disappear into the world.

Under Possibility~C, this is the cleanest explanation I have found for why Healer was simultaneously rejected by the intelligence community and protected by it. Under Possibilities~A and~B, it is pattern-matching applied to a famous case. The reader should weigh it accordingly.

\chapterreturn
